%%% ВЫБОР УСТРОЙСТВА %%%

\IfStrEqCase{\devicecode}{%
{ЮТКБ.656122.611-04.01 РЭ}{%
% код для этого типа
\def\Num{12}
}%
{ЮТКБ.656122.611-04.02 РЭ}{%
\def\Num{24}
}%
}
[
\def\Num{??}
]


\color{unimain}{\subsection{Логика отключения выключателей присоединений}}
\color{black}

\begin{enumerate}

\item
Ввод и вывод логики отключения выключателей присоединений осуществляется либо посредством программного ключа \textbf{«ЛО Вn»}, либо с использованием внешнего сигнала \textit{«Вывод действия на Bn»}.

\item
Сигнал отключения присоединения (\textit{«Отключение Вn»}) формируется при отключении СШ с учетом текущей фиксации присоединения, а также при отключении в результате опробования присоединения.

\item
Сигнал аварийного отключения выключателя присоединения (\textit{«Откл. аварийно Вn»}) дополнительно формируется при действии УРОВ на себя.

\item
Для каждого выключателя присоединения предусмотрена возможность формирования как общего, так и избирательного запрета АПВ.

\item
Общий запрет АПВ выключателя присоединения формируется по факту наличия сигнала запрет АПВ СШ (\textit{«Запрет АПВ 1С»} или \textit{«Запрет АПВ 2С»}) с учетом текущей фиксации присоединений. 

\item
Избирательный запрет АПВ выключателя присоединения формируется по факту срабатывания ДЗШ с учетом текущей фиксации присоединения.

\item 
Уставки ЛО выключателей приведены в таблицах \ref{app_set:lov1}-\ref{app_set:lov\Num}.

\item 
Алгоритм работы ЛО В1 приведен на рисунке \ref{pict:Log_otkl_Q1}, алгоритмы работы остальных выключателей выполнены аналогично.

\medskip
\begin{figure}[H]
\centering
\includegraphics[width=0.8\textwidth,height=0.8\textheight,keepaspectratio]{\fbpath/04.01 ДЗШ/761.9901 Логика отключения присоединений/img/Log_otkl_Q1.pdf}
\captionof{figure}{Функциональная схема логики отключения выключателя В1}\label{pict:Log_otkl_Q1}
\end{figure}


\end{enumerate}